\sect{12. Consumibles, varitas y anillos}

\textbf{Consultar un pergamino sin gastarlo.} Cógelo con el botón izquierdo, desde el suelo o desde una casilla del inventario: la zona de mensajes muestra «Pergamino: nombre del conjuro». Funciona con pergaminos de mago y divinos, incluso si el mago ya conoce el conjuro. Devuélvelo a una casilla con otro clic izquierdo. El botón derecho sobre la mano lo utiliza y consume. El nombre permite buscar su efecto en el catálogo de conjuros de este manual.

\textbf{Quién recibe el efecto.} Una poción la bebe quien la sostiene. Un pergamino o una varita que actúan sobre un personaje se aplican al portador al utilizarlos; pásalos antes al destinatario. Los conjuros preparados que piden un destinatario permiten escoger su retrato. No confundas ambos procedimientos.

\textbf{Un uso tiene un coste.} La comida, las pociones y los pergaminos desaparecen al consumirse. Las varitas pierden cargas y desaparecen al agotarse, salvo Silvias, que conserva sus cargas. No se recargan descansando. Un intento que exige una mano libre puede ser rechazado hasta que la liberes.

Al usar un pergamino o una varita, espera a que la mano vuelva a estar disponible.

\textbf{Identificación y duración.} Coge el objeto con el botón izquierdo y lanza Leer magia para identificarlo. Las pociones, varitas y anillos identificados muestran su efecto junto al nombre. Las nubes y los efectos temporales se conservan al guardar y cargar; el descanso hace avanzar su duración. Algunos efectos terminan al descansar, pero los de mayor duración pueden continuar.

\begin{reference}{Nombre}{Efecto, uso y límites}
\itemname{rations}{Raciones / raciones de hierro} & Coge la comida y déjala sobre el plato del inventario del personaje. Se consume y aumenta FOOD hasta un máximo de 100. Hay reservas de distinto tamaño; no todas alimentan lo mismo.\\[4pt]
\itemname{egg}{Huevo de kenku} & También se come sobre el plato. Su aportación de comida varía según el huevo. Conserva provisiones para los descansos largos.\\[4pt]
\itemname{potion-strength}{Poción de fuerza gigante} & Fija la fuerza en 22 durante aproximadamente 150--240 segundos. Después recupera la fuerza anterior, incluso si has bebido otra poción durante el efecto. No cura PV ni veneno.\\[4pt]
\itemname{potion-healing}{Poción de curación} & Restaura 2d4 + 2 PV (4--10), sin superar el máximo. La bebe el personaje que la sostiene.\\[4pt]
\itemname{potion-extra}{Poción de gran curación} & Restaura 3d8 + 3 PV (6--27), sin superar el máximo. No revive a los muertos.\\[4pt]
\itemname{potion-poison}{Poción de veneno} & Puede envenenar al bebedor si no resiste su efecto. Una botella sin identificar no es necesariamente una cura: guarda antes de experimentar.\\[4pt]
\itemname{potion-vitality}{Poción de vitalidad} & Llena FOOD hasta 100. Sacia al personaje, pero no es una poción de curación de PV.\\[4pt]
\itemname{potion-speed}{Poción de velocidad} & Permite usar las armas el doble de rápido durante unos minutos.\\[4pt]
\itemname{potion-invisible}{Poción de invisibilidad} & Aplica invisibilidad al bebedor: los monstruos tienen más dificultad para acertarle. No tiene un plazo fijo; usar una mano o lanzar magia ofensiva puede romperla, igual que el conjuro.\\[4pt]
\itemname{potion-antidote}{Antídoto} & Elimina el envenenamiento. Cura aparte el daño ya sufrido. La poción curativa enana también cura el veneno.\\[4pt]
\itemname{scroll-mage}{Pergamino de mago} & Úsalo desde una mano para gastar su conjuro una vez, o cópialo en el libro mediante Utilizar Pergaminos. Copiarlo también consume el pergamino.\\[4pt]
\itemname{scroll-divine}{Pergamino divino} & Contiene una oración de un solo uso. No se copia en el libro del mago. Si su efecto es sobre un personaje, el uso desde un objeto se aplica a quien lo sostiene.\\[4pt]
\end{reference}

\refspace
\subsection*{Varitas: efectos disponibles}

\begin{reference}{Nombre}{Efecto, uso y límites}
\itemname{wand}{Varita sin poder} & No produce un efecto mágico; aparece el aviso correspondiente.\\[4pt]
\itemname{wand}{Relámpagos} & Lanza Rayo eléctrico. Cada uso válido consume una carga.\\[4pt]
\itemname{wand}{Congelación} & Lanza Cono helado hacia delante. Orienta al grupo antes de usarla.\\[4pt]
\itemname{wand}{Curación} & Usa Curación grave (2d8 + 1 PV) sobre el portador. Pásala al personaje herido antes de activarla.\\[4pt]
\itemname{wand}{Bolas de fuego} & Lanza Bola de fuego. Necesita una trayectoria hacia el enemigo.\\[4pt]
\itemname{wand}{Silvias / Slivias} & Empuja al monstruo de delante una casilla por un paso abierto hasta una casilla vacía. Conserva sus cargas.\\[4pt]
\itemname{wand}{Misiles mágicos} & Lanza Dardo mágico. Es una reserva de ataques cuando faltan conjuros preparados.\\[4pt]
\end{reference}

\refspace
\subsection*{Anillos: llevarlos puestos}

\begin{reference}{Nombre}{Efecto, uso y límites}
\itemname{ring-protection}{Protección} & Su bonificación mejora la CA según las reglas del equipo de la sección 11. Llevarlo en la mochila no basta.\\[4pt]
\itemname{ring}{Sustento} & Equipado en un dedo, evita la pérdida normal de comida, también durante el descanso. No cura heridas.\\[4pt]
\itemname{ring}{Ornamento} & Adorno. No da bonificaciones de combate.\\[4pt]
\itemname{ring}{Hechicería} & Equipado en un dedo, duplica los huecos de preparación de mago para conjuros de nivel 4 y 5. Llevar dos no multiplica de nuevo la bonificación. No concede acceso a niveles superiores ni conocimiento de nuevos hechizos.\\[4pt]
\itemname{ring}{Caída suave} & Equipado en un dedo, evita el daño de caída al que se aplica esta protección. No protege frente a todas las trampas ni frente a los ataques normales.\\[4pt]
\end{reference}
